\chapter{Комбинационные схемы}
\label{ch:comb}

В главе~\ref{ch:gates} мы научились записывать логические функции и собирать
их из элементов, но каждый раз функция была дана готовой. На практике
задача звучит иначе: «тревога, если открыта дверь и сейчас ночь или включена
охрана». Как превратить такую фразу в схему и не сделать лишних
элементов? И второй вопрос, оставшийся с прошлой главы: мы считали элементы
мгновенными, но настоящей схеме нужно время, чтобы «успокоиться». Что
происходит с выходом в этот момент? Эта глава отвечает на оба вопроса.

\begin{buildbox}
Откройте вкладку \emph{Обучение}, книгу «Логические элементы», урок
\emph{«Два из трёх»}: соберите схему, в которой лампа $y$ горит, если из трёх
кнопок $a$, $b$, $c$ нажаты хотя бы две. Не читая дальше, запишите в тетрадь,
сколько блоков вы использовали. В конце главы мы получим эту же схему
по плану «таблица $\to$ выражение $\to$ упрощение $\to$ схема» и сравним.
\end{buildbox}

\section{Комбинационная и последовательностная логика}

Все цифровые схемы делятся на два больших класса.

\begin{description}
  \item[Комбинационные схемы] --- выходы зависят \emph{только от текущих
    значений входов}. У них нет памяти: подайте те же входы --- получите те же
    выходы, неважно, что было раньше. Примеры: логические элементы,
    дешифраторы, мультиплексоры, сумматоры.
  \item[Последовательностные схемы] --- выходы зависят не только от текущих
    входов, но и от \emph{предыстории}, то есть от того, что было раньше.
    У них есть память. Примеры: триггеры, регистры, счётчики. До них мы
    дойдём позже; пока достаточно знать, что они существуют и что
    «запомнить, что было» комбинационная схема не умеет.
\end{description}

Признак, по которому их легко различить на рисунке: у комбинационной схемы
сигнал идёт строго от входов к выходам, а у последовательностной часть
выходов возвращается обратно на вход через память.

\begin{figure}[H]
\centering
\begin{tikzpicture}
  \node[block, minimum width=30mm, minimum height=20mm] (c) {Комбинационная\\логика};
  \foreach \i in {1,2,3} \draw[arr] ($(c.west)+(-1.2,0.9-\i*0.45)$) -- ($(c.west)+(0,0.9-\i*0.45)$);
  \foreach \i in {1,2} \draw[arr] ($(c.east)+(0,0.68-\i*0.45)$) -- ($(c.east)+(1.2,0.68-\i*0.45)$);
  \node[lbl, left=12mm of c] {входы};
  \node[lbl, right=12mm of c] {выходы};
  \node[font=\sffamily\small\bfseries, above=2mm of c] {Комбинационная схема};
  \begin{scope}[xshift=8.3cm]
    \node[block, minimum width=30mm, minimum height=20mm] (s) {Комбинационная\\логика};
    \node[gblock, minimum width=30mm, below=6mm of s] (m) {Память\\(триггеры)};
    \foreach \i in {1,2,3} \draw[arr] ($(s.west)+(-1.2,0.9-\i*0.45)$) -- ($(s.west)+(0,0.9-\i*0.45)$);
    \foreach \i in {1,2} \draw[arr] ($(s.east)+(0,0.68-\i*0.45)$) -- ($(s.east)+(1.2,0.68-\i*0.45)$);
    \draw[arr] ($(s.east)+(0.4,-0.7)$) -- ++(0,-1.0) |- (m.east);
    \draw[arr] (m.west) -- ++(-0.4,0) |- ($(s.west)+(0,-0.7)$);
    \node[lbl, left=4mm of m, align=center] {состояние};
    \node[font=\sffamily\small\bfseries, above=2mm of s] {Последовательностная схема};
  \end{scope}
\end{tikzpicture}
\caption{Главное отличие: у последовательностной схемы есть память и обратная
связь через неё}
\label{fig:comb_vs_seq}
\end{figure}

\begin{diybox}[Есть ли у схемы память]
\textbf{Задание.} \stars{1} Для каждой схемы реши, комбинационная она или нет:
лампа «горит, пока нажата кнопка»; лампа, которая «включается первым
нажатием и выключается вторым»; сигнал «сумма двух чисел».
\stars{2} Придумай бытовое устройство каждого класса.

\textbf{В Logra.} Вкладка \emph{Конструктор}: соберите «кнопка $\to$ НЕ $\to$
лампа» и подайте на кнопку 0, потом 1, потом снова 0. Дальше откройте
пример \code{edge-toggle} (библиотека примеров, вкладка \emph{Код}) и
удерживайте кнопку несколько тактов.

\textbf{Ожидаемый результат.} В первой схеме одни и те же входы всегда дают
одни и те же выходы. Во второй переключение происходит один раз
за нажатие: схема «помнит», что кнопка уже была нажата.

\textbf{Как проверить.} Вкладка \emph{Диаграммы}: у комбинационной схемы
выход можно предсказать по таблице истинности, глядя на текущие входы;
у \code{edge-toggle} --- нельзя, нужна ещё история.
\end{diybox}

\section{От словесного условия к схеме}

Разработка комбинационной схемы всегда идёт по одному плану:
\begin{enumerate}
  \item понять, какие входы и выходы нужны, и дать им имена;
  \item составить \textbf{таблицу истинности};
  \item записать логическое выражение;
  \item \textbf{упростить} его;
  \item нарисовать схему.
\end{enumerate}

Главная трудность --- первый шаг. Фраза на русском языке почти никогда не
говорит, сколько входов у задачи и что делать в «пограничных» случаях.
Поэтому перед таблицей стоит выписать для каждого входа, что значит 1 и что
значит 0 («дверь открыта: 1», «ночь: 1»), и только потом строить таблицу.

\subsection{Пример: голосование трёх судей}

На соревнованиях работают три судьи: $a$, $b$, $c$. Каждый нажимает кнопку,
если засчитывает попытку. Лампа $y$ должна загораться, если попытку
засчитали \emph{хотя бы двое}. Такая функция называется
\term{мажоритарной} (от лат. \emph{major} --- большинство).

\textbf{Шаг 1--2. Таблица истинности.} Три входа --- $2^3 = 8$ строк:

\begin{table}[H]
\centering
\caption{Таблица истинности мажоритарной функции}
\label{tab:maj}
\begin{tabular}{c|ccc|c|l}
\toprule
\textbf{№} & $a$ & $b$ & $c$ & $y$ & \textbf{Слагаемое} \\
\midrule
0 & 0 & 0 & 0 & 0 & \\
1 & 0 & 0 & 1 & 0 & \\
2 & 0 & 1 & 0 & 0 & \\
3 & 0 & 1 & 1 & \textbf{1} & $\overline{a}bc$ \\
4 & 1 & 0 & 0 & 0 & \\
5 & 1 & 0 & 1 & \textbf{1} & $a\overline{b}c$ \\
6 & 1 & 1 & 0 & \textbf{1} & $ab\overline{c}$ \\
7 & 1 & 1 & 1 & \textbf{1} & $abc$ \\
\bottomrule
\end{tabular}
\end{table}

\textbf{Шаг 3. Выражение.} Для каждой строки, где $y=1$, запишем
произведение всех входов: вход берётся как есть, если он равен 1, и с
инверсией, если равен 0. Такое произведение равно единице \emph{только} в
этой строке. Сложив все такие произведения, получим функцию. Эта форма
записи называется \term{совершенной дизъюнктивной нормальной формой} (СДНФ):
\[
  y = \overline{a}bc + a\overline{b}c + ab\overline{c} + abc.
\]

\textbf{Шаг 4. Упрощение.} Слагаемое $abc$ можно повторить сколько угодно
раз ($x + x = x$). Добавим его ещё дважды и сгруппируем:
\[
  y = (\overline{a}bc + abc) + (a\overline{b}c + abc) + (ab\overline{c} + abc)
    = bc(\overline{a}+a) + ac(\overline{b}+b) + ab(\overline{c}+c)
    = bc + ac + ab.
\]
Вместо четырёх трёхвходовых элементов И и трёх инверторов нужно всего три
двухвходовых элемента И и один элемент ИЛИ.

\textbf{Шаг 5. Схема.}

\begin{figure}[H]
\centering
\begin{tikzpicture}
  \foreach \n/\x in {a/0, b/0.6, c/1.2} {
    \draw[very thick, FpgaBlue] (\x,0.6) node[above, text=FpgaDark] {$\n$} -- (\x,-3.4);
  }
  \node[gAND] (g1) at (3.4,0) {};
  \node[gAND] (g2) at (3.4,-1.4) {};
  \node[gAND] (g3) at (3.4,-2.8) {};
  \node[gOR, logic gate inputs=nnn, minimum height=14mm] (o) at (6,-1.4) {};
  \draw[wire] (g1.input 1) -- (0,0 |- g1.input 1) node[circle, fill, inner sep=1.2pt] {};
  \draw[wire] (g1.input 2) -- (0.6,0 |- g1.input 2) node[circle, fill, inner sep=1.2pt] {};
  \draw[wire] (g2.input 1) -- (0,0 |- g2.input 1) node[circle, fill, inner sep=1.2pt] {};
  \draw[wire] (g2.input 2) -- (1.2,0 |- g2.input 2) node[circle, fill, inner sep=1.2pt] {};
  \draw[wire] (g3.input 1) -- (0.6,0 |- g3.input 1) node[circle, fill, inner sep=1.2pt] {};
  \draw[wire] (g3.input 2) -- (1.2,0 |- g3.input 2) node[circle, fill, inner sep=1.2pt] {};
  \draw[wire] (g1.output) -- node[above, lbl] {$ab$} ++(0.6,0) |- (o.input 1);
  \draw[wire] (g2.output) -- node[above, lbl] {$ac$} (o.input 2);
  \draw[wire] (g3.output) -- node[below, lbl] {$bc$} ++(0.6,0) |- (o.input 3);
  \draw[wire] (o.output) -- ++(0.8,0) node[right] {$y$};
\end{tikzpicture}
\caption{Мажоритарный элемент «2 из 3»: $y = ab + ac + bc$}
\label{fig:maj}
\end{figure}

\subsection{Пример: охрана с ночным режимом}

Теперь словесная задача, которая нам пригодится дальше. Сигнализация
должна звонить, если дверь открыта \emph{и} (сейчас ночь \emph{или}
включена охрана). Назовём входы: $d$ --- дверь открыта, $n$ --- ночь,
$s$ --- охрана включена. Скобки в условии показывают порядок: сначала
«ночь или охрана», потом «и дверь»:
\[
  y = d\cdot(n + s).
\]
Если раскрыть скобки, получится $y = dn + ds$: два элемента И и один ИЛИ
против одного ИЛИ и одного И в исходной записи. Скобочная форма
\emph{экономнее} (два элемента вместо трёх) и по смыслу честнее: она
повторяет условие слово в слово. Как выбирать между такими формами, мы
подробно разберём в следующей главе.

\begin{diybox}[Из слов в схему]
\textbf{Задание.} \stars{1} Соберите схему «2 из 3» по рисунку~\ref{fig:maj}
и проверьте все 8 комбинаций. \stars{2} Измените условие: лампа должна
гореть, если засчитали \emph{ровно двое}. Составьте таблицу и выражение.
\stars{3} Этап проекта «Охранная панель»: добавьте к «2 из 3» вход
охраны $arm$: тревога только при $arm=1$ и согласии хотя бы двух из трёх
датчиков двери.

\textbf{В Logra.} \stars{1} Вкладка \emph{Обучение}: урок
\emph{«Два из трёх»} (книга «Логические элементы»); пример
\code{majority}. \stars{2} Пример \code{exactly-two}: сравните со своим
выражением. \stars{3} Вкладка \emph{Обучение}, книга «Проекты»,
урок \emph{«Голосование «2 из 3» с охраной»}.

\textbf{Ожидаемый результат.} В \stars{1} совпадают все 8 строк таблицы
\ref{tab:maj}. В \stars{2} для $111$ лампа не горит. В \stars{3}
при $arm=0$ лампа не горит ни при каких $a,b,c$; единица на выходе
стоит ровно в четырёх строках из шестнадцати ($arm\,abc=1011$, $1101$,
$1110$, $1111$).

\textbf{Как проверить.} Вкладка \emph{Диаграммы}: задайте $a,b,c$ по
очереди как счётчик 0--7 и сверьте $y$ с таблицей; урок проверяется
встроенным симулятором (кнопка проверки). Критерий приёмки этапа~3:
четыре строки с $y=1$ из шестнадцати и тревога погашена при $arm=0$.
Минимальную форму этой функции мы найдём в главе~\ref{ch:minimize}.
\end{diybox}

\section{Карты Карно}

Упрощать выражения алгебраически можно, но легко ошибиться. Для функций от
2--4 переменных есть наглядный способ --- \term{карта Карно}. Это та же
таблица истинности, но нарисованная в виде прямоугольника так, что
\emph{соседние клетки отличаются значением только одной переменной}.
Поэтому столбцы подписаны в порядке 00, 01, 11, 10 (код Грея), а не
00, 01, 10, 11.

Правило упрощения: единицы объединяют в прямоугольные группы по 2, 4, 8
клеток (степень двойки), стараясь сделать группы как можно больше. Каждой
группе соответствует одно произведение, в котором остаются только те
переменные, что \emph{не меняются} внутри группы.

\begin{figure}[H]
\centering
\begin{tikzpicture}[cell/.style={draw, minimum size=11mm, font=\sffamily\large}]
  \foreach \c/\lab in {0/00, 1/01, 2/11, 3/10} \node[font=\sffamily\small] at (\c*1.1,1.05) {\lab};
  \foreach \r/\lab in {0/0, 1/1} \node[font=\sffamily\small] at (-0.95,-\r*1.1) {\lab};
  \node[font=\sffamily\small] at (-1.1,1.05) {$a \backslash bc$};
  % значения: строка a=0: bc=00,01,11,10 -> 0,0,1,0; a=1: 0,1,1,1
  \foreach \r/\c/\v in {0/0/0, 0/1/0, 0/2/1, 0/3/0, 1/0/0, 1/1/1, 1/2/1, 1/3/1}
    {\ifnum\v=1 \def\cc{FpgaLight}\else\def\cc{white}\fi
     \node[cell, fill=\cc] at (\c*1.1,-\r*1.1) {\v};}
  \draw[very thick, FpgaRed, rounded corners=4pt] (1.7,-1.62) rectangle (2.7,0.52);
  \draw[very thick, FpgaGreen, rounded corners=4pt] (0.6,-1.56) rectangle (2.62,-0.64);
  \draw[very thick, FpgaOrange, rounded corners=4pt] (1.78,-1.48) rectangle (3.8,-0.72);
  \node[font=\sffamily\small, text=FpgaRed, anchor=west] at (4.6,0.2) {$bc$ --- столбец $bc=11$ ($a$ меняется)};
  \node[font=\sffamily\small, text=FpgaGreen, anchor=west] at (4.6,-0.5) {$ac$ --- клетки $abc=101, 111$};
  \node[font=\sffamily\small, text=FpgaOrange, anchor=west] at (4.6,-1.2) {$ab$ --- клетки $abc=111, 110$};
  \node[font=\sffamily\small\bfseries, anchor=west] at (4.6,-2.0) {$y = ab + ac + bc$};
\end{tikzpicture}
\caption{Карта Карно мажоритарной функции: три группы по две клетки дают
три слагаемых}
\label{fig:karnaugh}
\end{figure}

Разберём, откуда берутся слагаемые. Красная группа --- столбец $bc=11$: в
двух его клетках меняется только $a$, а $b=1$ и $c=1$ остаются, значит
слагаемое $bc$. Зелёная группа --- клетки $101$ и $111$: тут меняется $b$,
остаются $a=1$, $c=1$, и слагаемое $ac$. Оранжевая: меняется $c$, остаются
$a=1$, $b=1$, слагаемое $ab$. Три группы дают те же три слагаемых, что мы
получили алгеброй, но без выкладок.

Две особенности трёхпеременной карты, о которых забывают. Во-первых, карта
\emph{замкнута по горизонтали}: крайние столбцы $00$ и $10$ тоже соседи
(отличаются только $b$). Во-вторых, одна и та же единица может входить в
несколько групп: пересечение групп разрешено и часто выгодно, оно делает
группы больше.

\begin{infobox}[Кто будет упрощать схемы в больших проектах?]
Когда схемы становятся большими, упрощение выполняет специальная программа,
причём гораздо лучше человека. Но понимать, \emph{как} это делается, всё
равно полезно: это помогает оценить, сколько ресурсов займёт ваша схема и
насколько быстро она будет работать. Полный набор приёмов (четыре
переменные, безразличные состояния, формальный алгоритм) --- в
следующей главе.
\end{infobox}

\begin{diybox}[Карта Карно на бумаге и в схеме]
\textbf{Задание.} \stars{1} Для $F_2=\Sigma m(1,3,5,7)$ нарисуйте карту,
найдите минимальную форму и соберите её в Конструкторе. \stars{2} Для
$F_3=\Sigma m(3,5,6,7)$ найдите минимальную форму, соберите и сравните с
схемой «2 из 3» (это та же функция в другой записи?). \stars{3} Придумайте
функцию трёх переменных, у которой минимальная форма состоит из трёх
слагаемых по две буквы, но группы \emph{не} образуют ни одного «замыкания
через край».

\textbf{В Logra.} Вкладка \emph{Конструктор}: схема из блоков
И, ИЛИ, НЕ. Входы $a,b,c$ задайте переключателями.

\textbf{Ожидаемый результат.} Для $F_2$ минимальная форма --- один
проводок, $F_2=c$; в схеме нет ни одного элемента между входом и выходом.
Для $F_3$ получится $ab+ac+bc$, то есть та же мажоритарная функция.

\textbf{Как проверить.} Вкладка \emph{Диаграммы}: переберите все 8
комбинаций; выход должен совпасть со столбцом $F$ вашей карты. Вкладка
\emph{RTL}: сравните свою схему с рисунком из этой главы.
\end{diybox}

\section{Задержка распространения}

До сих пор мы считали, что элемент срабатывает мгновенно. На самом деле
транзисторам нужно время, чтобы переключиться, а проводам --- чтобы
зарядиться. Выход элемента меняется через время $t_\text{зд}$ после входа.
Это \term{задержка распространения}. У современных микросхем она составляет
от долей наносекунды до нескольких наносекунд.

\begin{figure}[H]
\centering
\begin{tikzpicture}[x=1cm, y=1cm]
  \draw[very thick, FpgaBlue] (0,1.6) -- (2,1.6) -- (2.1,2.4) -- (6,2.4) -- (6.1,1.6) -- (9,1.6);
  \draw[very thick, FpgaOrange] (0,0.8) -- (3,0.8) -- (3.1,0) -- (7,0) -- (7.1,0.8) -- (9,0.8);
  \node[font=\sffamily\small, anchor=east] at (-0.1,2.0) {$a$};
  \node[font=\sffamily\small, anchor=east] at (-0.1,0.4) {$\overline{a}$};
  \draw[dashed, FpgaGray] (2.05,2.6) -- (2.05,-0.4); \draw[dashed, FpgaGray] (3.05,2.6) -- (3.05,-0.4);
  \draw[<->, thick] (2.05,-0.3) -- (3.05,-0.3) node[midway, below, font=\sffamily\small] {$t_\text{зд}$};
  \draw[dashed, FpgaGray] (6.05,2.6) -- (6.05,-0.4); \draw[dashed, FpgaGray] (7.05,2.6) -- (7.05,-0.4);
  \draw[<->, thick] (6.05,-0.3) -- (7.05,-0.3) node[midway, below, font=\sffamily\small] {$t_\text{зд}$};
\end{tikzpicture}
\caption{Выход инвертора меняется с задержкой $t_\text{зд}$}
\end{figure}

Если сигнал проходит через несколько элементов подряд, задержки
складываются. Самый длинный путь от входа до выхода называется
\term{критическим путём}. Он определяет, насколько быстро может работать
схема: пока на выходе не установилось правильное значение, результат
читать нельзя.

\begin{figure}[H]
\centering
\begin{tikzpicture}
  \node[gAND, minimum width=10mm, minimum height=7mm] (g1) at (0,0) {};
  \node[gOR, minimum width=10mm, minimum height=7mm] (g2) at (1.8,-0.3) {};
  \node[gXOR, minimum width=10mm, minimum height=7mm] (g3) at (3.6,-0.6) {};
  \node[gNOT] (g4) at (5.3,-0.6) {};
  \draw[wire, FpgaRed, very thick] (g1.input 1) -- ++(-0.5,0) node[left, text=FpgaDark] {$a$};
  \draw[wire] (g1.input 2) -- ++(-0.5,0) node[left] {$b$};
  \draw[wire, FpgaRed, very thick] (g1.output) -- ++(0.2,0) |- (g2.input 1);
  \draw[wire] (g2.input 2) -- ++(-2.3,0) node[left] {$c$};
  \draw[wire, FpgaRed, very thick] (g2.output) -- ++(0.2,0) |- (g3.input 1);
  \draw[wire] (g3.input 2) -- ++(-4.1,0) node[left] {$d$};
  \draw[wire, FpgaRed, very thick] (g3.output) -- (g4.input);
  \draw[wire, FpgaRed, very thick] (g4.output) -- ++(0.6,0) node[right, text=FpgaDark] {$y$};
  \node[font=\sffamily\small, text=FpgaRed, anchor=west] at (7.2,-0.2) {критический путь:};
  \node[font=\sffamily\small, text=FpgaRed, anchor=west] at (7.2,-0.8) {$t = 4\,t_\text{зд}$};
\end{tikzpicture}
\caption{Критический путь проходит через четыре элемента}
\end{figure}

Простое следствие: выигрыш от упрощения схемы бывает не только в числе
элементов, но и во времени. Если вместо цепочки из четырёх элементов
удаётся поставить два, критический путь сокращается вдвое.

\subsection{Состязания и «иголки»}

Из-за задержек на выходе комбинационной схемы могут появляться короткие
ложные импульсы --- \term{иголки} (glitches). Классический пример:
$y = a\cdot\overline{a}$. По законам алгебры $y$ всегда равен 0. Но инвертор
срабатывает с задержкой, и в течение короткого времени после фронта $a$
\emph{оба} входа элемента И равны 1.

\begin{figure}[H]
\centering
\begin{tikzpicture}
  \node[gAND, minimum width=10mm, minimum height=7mm] (g) at (2.2,0) {};
  \node[gNOT] (n) at (0.9,-0.7) {};
  \draw[wire] (-0.6,0.12) node[left] {$a$} -- (g.input 1);
  \fill (0.0,0.12) circle (1.3pt);
  \draw[wire] (0.0,0.12) |- (n.input);
  \draw[wire] (n.output) -- ++(0.3,0) |- (g.input 2);
  \draw[wire] (g.output) -- ++(0.5,0) node[right] {$y$};
  \begin{scope}[xshift=5cm, yshift=0.9cm, x=1cm, y=1cm]
    \draw[very thick, FpgaBlue] (0,0) -- (1.5,0) -- (1.55,0.5) -- (5,0.5);
    \draw[very thick, FpgaOrange] (0,-0.4) -- (2.4,-0.4) -- (2.45,-0.9) -- (5,-0.9);
    \draw[very thick, FpgaRed] (0,-1.8) -- (1.6,-1.8) -- (1.65,-1.3) -- (2.5,-1.3) -- (2.55,-1.8) -- (5,-1.8);
    \node[font=\sffamily\small, anchor=east] at (-0.1,0.25) {$a$};
    \node[font=\sffamily\small, anchor=east] at (-0.1,-0.65) {$\overline{a}$};
    \node[font=\sffamily\small, anchor=east] at (-0.1,-1.55) {$y$};
    \node[font=\sffamily\small, text=FpgaRed, anchor=west] at (2.7,-1.4) {ложный импульс!};
  \end{scope}
\end{tikzpicture}
\caption{Возникновение «иголки» из-за задержки инвертора}
\end{figure}

Это пример \emph{состязания}: два пути от одного входа к выходу (прямой и
через инвертор) имеют разную длину, и сигналы по ним «бегут наперегонки».
Выход на мгновение принимает значение, которого нет в таблице истинности.

\subsubsection{Иголка и карта Карно}

Бывает, что иголка возникает даже в схеме, где логика безупречна. Возьмём
$F = AB + \overline{A}C$ и зафиксируем $B=C=1$. Тогда $F = A + \overline{A}$
всегда должна быть равна 1. Но слагаемое $AB$ гаснет, когда $A$ переходит из
1 в 0, а слагаемое $\overline{A}C$ загорается позже, потому что сигнал
$\overline{A}$ ещё «в пути» через инвертор. На короткое время гаснут
\emph{оба} слагаемых, и на выходе $F$ появляется нулевая иголка.

На карте Карно это видно сразу: две группы $AB$ и $\overline{A}C$
\emph{соприкасаются}, но не перекрываются. Переход между ними идёт по
границе, где ни одна группа не удерживает единицу. Лекарство --- добавить
перекрывающую группу $BC$: она накрывает обе клетки на границе и держит
выход в единице, пока $A$ переключается.

\begin{figure}[H]
\centering
\begin{tikzpicture}[cell/.style={draw, minimum size=11mm, font=\sffamily\large}]
  \foreach \c/\lab in {0/00, 1/01, 2/11, 3/10} \node[font=\sffamily\small] at (\c*1.1,1.05) {\lab};
  \foreach \r/\lab in {0/0, 1/1} \node[font=\sffamily\small] at (-0.95,-\r*1.1) {\lab};
  \node[font=\sffamily\small] at (-1.1,1.05) {$A \backslash BC$};
  \foreach \r/\c/\v in {0/0/0, 0/1/1, 0/2/1, 0/3/0, 1/0/0, 1/1/0, 1/2/1, 1/3/1}
    {\ifnum\v=1 \def\cc{FpgaLight}\else\def\cc{white}\fi
     \node[cell, fill=\cc] at (\c*1.1,-\r*1.1) {\v};}
  \draw[very thick, FpgaOrange, rounded corners=4pt] (0.65,0.42) rectangle (2.55,-0.42);
  \draw[very thick, FpgaGreen, rounded corners=4pt] (1.85,-0.68) rectangle (3.75,-1.52);
  \draw[very thick, dashed, FpgaRed, rounded corners=4pt] (1.9,0.5) rectangle (2.5,-1.6);
  \node[font=\sffamily\small, text=FpgaOrange, anchor=west] at (4.6,0.2) {$\overline{A}C$ --- клетки $A=0$};
  \node[font=\sffamily\small, text=FpgaGreen, anchor=west] at (4.6,-0.5) {$AB$ --- клетки $A=1$};
  \node[font=\sffamily\small, text=FpgaRed, anchor=west] at (4.6,-1.2) {$BC$ --- добавка против иголки};
  \node[font=\sffamily\small\bfseries, anchor=west] at (4.6,-2.0) {$F = AB + \overline{A}C + BC$};
\end{tikzpicture}
\caption{Группы $\overline{A}C$ и $AB$ соприкасаются, но не перекрываются.
Пунктирная группа $BC$ устраняет иголку, хотя и не нужна по логике}
\label{fig:hazard_karnaugh}
\end{figure}

Парадокс: добавленная группа \emph{лишняя} с точки зрения минимальной
формы, но нужна для надёжной работы. Поэтому «минимальная» и «лучшая» схемы
бывают разными, и это нужно помнить.

\begin{warnbox}[Как с этим бороться]
Иголки --- главная причина, по которой в серьёзных схемах выходы
комбинационной логики «читают» не когда попало, а только по фронту
тактового сигнала, когда все переходные процессы уже закончились. Это
называется \term{синхронным проектированием} (глава~\ref{ch:sync}).
Практическое правило: иголка безвредна, если её никто не читает, пока она
не закончилась.
\end{warnbox}

\begin{diybox}[Задержка и иголка]
\textbf{Задание.} \stars{1} Для схемы с рисунка «Критический путь»
посчитайте задержку, если $t_\text{зд}=1{,}2$~нс у каждого элемента, и
нарисуйте диаграмму $y=a\cdot\overline a$ с задержкой инвертора, отметив
иголку. \stars{2} Для $F=AB+\overline{A}C$ при $B=C=1$ нарисуйте диаграмму
и найдите момент, когда $F$ на мгновение падает в 0. Добавьте слагаемое
$BC$ и повторите. \stars{3} Придумайте схему, в которой иголка вредна
(подсказка: выход подключён к входу, реагирующему на любой короткий
импульс).

\textbf{В Logra.} Вкладка \emph{Конструктор}: соберите $y=a\cdot\overline a$
(вход $a$, блок НЕ, блок И), а рядом цепочку из пяти блоков НЕ. Затем
пример \code{button-debounce} (библиотека примеров): как дребезг контакта
похож на иголку.

\textbf{Ожидаемый результат.} Для расчёта: $4\cdot 1{,}2 = 4{,}8$~нс. Пять
инверсий дают $y=\overline a$. На диаграмме выход $y=a\cdot\overline a$
остаётся нулём при любом $a$: обычная симуляция логики не моделирует
задержки элементов. Это не ошибка, а важное наблюдение: иголки на такой
диаграмме не видно, а в реальной схеме она есть.

\textbf{Как проверить.} Вкладка \emph{Диаграммы}: убедитесь, что $y\equiv0$
(логика верна), и сравните с вашей бумажной диаграммой с иголкой. Разница
между ними --- в точности та проблема, ради которой в схемах с памятью
вводят общий тактовый сигнал. Откуда берутся наносекунды --- в
главе~\ref{ch:electrical}.
\end{diybox}

\begin{warnbox}[Типичные ошибки]
\begin{itemize}
  \item \textbf{Упрощать «как получится».} Алгебраические преобразования
        легко испортить: пропустить инверсию, забыть, что $\overline{a+b}\neq\overline a+\overline b$.
        После каждого упрощения сверяйте результат с таблицей истинности
        хотя бы на трёх строках.
  \item \textbf{Брать в группу не те клетки.} Группа должна быть
        прямоугольником из $1,2,4,8$ клеток. Три клетки подряд, «уголок»
        и диагональ --- не группа.
  \item \textbf{Забыть про соседство по краю.} Столбцы $00$ и $10$ на
        карте --- соседи. Нашли две единицы по краям и не объединили --- получили
        лишнее слагаемое.
  \item \textbf{Считать иголку ошибкой логики.} Иголка не нарушает
        таблицу истинности: она живёт в переходный момент и убирается
        тактированием или добавкой перекрывающей группы.
\end{itemize}
\end{warnbox}

\begin{ideabox}
Комбинационная схема не имеет памяти. Её проектируют по таблице истинности,
упрощают алгебраически или по карте Карно. Реальные элементы имеют
задержку, из-за которой возможны кратковременные ложные импульсы.
\end{ideabox}

\section{Вопросы и задачи}
\label{sec:tasks-comb}

Ответы собраны в приложении~\ref{app:answers}.

\begin{taskbox}[Задачи к главе]
\begin{enumerate}
  \item \stars{1}~Найдите по карте Карно минимальную форму функций $F_1=\Sigma m(0,1,4,5)$, $F_2=\Sigma m(1,3,5,7)$ и $F_3=\Sigma m(3,5,6,7)$ от трёх переменных $A,B,C$ ($A$ --- старший разряд).
  \item \stars{1}~Сигнал проходит через 5 элементов с задержкой 1,5~нс каждый. Чему равна задержка этого пути? Как она связана с предельной частотой схемы?
  \item \stars{2}~Сигнализация звонит, если дверь открыта \emph{и} (сейчас ночь \emph{или} включена охрана). Запишите функцию и нарисуйте схему.
  \item \stars{1}~Какие из узлов комбинационные, а какие последовательностные: сумматор, счётчик, мультиплексор, кнопка, которая переключает лампу при каждом нажатии?
  \item \stars{2}~У функции $F = AB + \overline{A}C$ при $B=C=1$ меняется $A$. Возникает ли «иголка»? Какое слагаемое нужно добавить, чтобы её убрать?
  \item \stars{2}~\emph{Найдите ошибку.} Студент упростил $y=\overline{a}bc+a\overline{b}c+ab\overline{c}+abc$ и получил $y=ab+c$. Проверьте его ответ по таблице истинности и найдите строку, на которой он ошибся.
  \item \stars{2}~\emph{Найдите ошибку.} На карте трёх переменных для $F=\Sigma m(1,2,3)$ ученик обвёл «уголок» из трёх клеток $1,3,2$ одной группой и записал $F=\overline{A}$. Чему равна его функция на самом деле? Какой должна быть минимальная форма?
  \item \stars{3}~\emph{Проект с критерием.} Построить схему «охрана с ночным режимом» $y=d\,(n+s)$ в Конструкторе и убедиться, что для всех 8 комбинаций $y$ совпадает с таблицей. Критерий приёмки: 8 из 8 строк совпали на вкладке Диаграммы.
  \item \stars{3}~\emph{Проект с критерием (этап проекта «Охранная панель»).} Добавить к схеме «2 из 3» вход охраны $arm$ и собрать минимальную форму $y=arm\cdot(ab+ac+bc)$. Критерий приёмки: при $arm=0$ выход равен нулю при любых $a,b,c$; при $arm=1$ выход совпадает с таблицей~\ref{tab:maj}. Сколько элементов получилось?
  \item \stars{3}~Покажите, что для функции $y=ab+\overline{a}c$ иголка возникает при $b=c=1$, а для функции $y=ab+\overline{a}c+bc$ --- нет. Как это связано с картой Карно?
\end{enumerate}
\end{taskbox}

\section{Что дальше}

Мы умеем строить схему по таблице и упрощать функцию от трёх переменных на
карте Карно. Но функций от четырёх переменных гораздо больше, у некоторых
комбинаций вообще нет значения, а реальные программы упрощают функции от
десятков переменных. В следующей главе мы разберём карту Карно на четыре
переменные, безразличные состояния, неверные группировки и посмотрим,
как минимальная форма превращается в настоящую схему в Logra.
